Calculs de l'auteur (synthese/synthese\_ensemble.py), d'après les sorties des quatre parties : sorties/proposition\_ircsg.json, controle\_hauts\_purs.txt, pfu/sorties/calculs\_pfu.json.\par Note : 1 \% des foyers au revenu le plus élevé ; convention centrale du modèle IR-CSG pour la répartition des revenus du capital entre tranches. Fin de l'effacement : moyenne des dix premières années (0,7 à 6,8 Md\euro{}), en supposant que le 1 \% en paie 90 \%. CEHR : 1,4 Md\euro{} (revenus de 2025), dont 50 à 80 \% payés sur les revenus du capital (hypothèse). L'immobilier (au plus quelques centaines de millions) et l'IS ne sont pas comptés. Losanges : même bilan avec 10,2 Md\euro{} d'IR imputés aux revenus du capital au lieu de 11,5 (revenus de 2024), ou dans la convention « haut ».
